% Source pp. 27-37 (Lectures1-6.pdf).
\section{More on Quantum Klein-Gordon Theory}

The space of states $\Ham$ of the quantum KG theory is the Fock space, defined
as the linear envelope of the basic states
\[
\ket{\pp_1,\dots,\pp_N}=\hat A^\dagger_{\pp_1}\cdots\hat A^\dagger_{\pp_N}\ket{0}\,,
\]
where the vacuum $\ket0$ satisfies
\[
\hat A_\pp\ket0=0\qquad\text{for all}\quad\pp\,.
\]
Here I have returned to the default case of the infinite Minkowski space-time,
so that the 3-momenta $\pp_i$ may take continuous unrestricted values.

The states $\ket{\pp_1,\dots,\pp_N}$ are interpreted as $N$-particle states.
Their scalar products follow from the commutation relations
$[\hat A_\pp,\hat A^\dagger_{\pp'}]=2\omega_\pp\,(2\pi)^3\,\delta^{(3)}(\pp-\pp')$.
In particular, for the one-particle states
\[
\braket{\pp}{\pp'}=2\omega_\pp\,(2\pi)^3\,\delta^{(3)}(\pp-\pp')\,.
\]
Since the operators $\hat A^\dagger_\pp$ are Lorentz invariant (only the
momentum $\pp$ transforms under the Lorentz transformations), the particles
have zero spin. Since the creation operators $\hat A^\dagger_\pp$ commute, the
$N$-particle states are symmetric with respect to permutations of the momenta
$\pp_i$: the particles are bosons.

The spectrum of the energy-momentum operator $(\hat H-E_0,\hat{\PPv})$ is the
composition of spectra of the $N$-particle states. In the space of vectors
$(E-E_0,\PPv)$ the vacuum resides at the origin of coordinates, the one
particle states form the mass-shell hyperboloid $(E-E_0)^2-\PPv^2=m^2$, the
two-particle states fill the inside of the hyperboloid
$(E-E_0)^2-\PPv^2=4m^2$, etc.\ (figure~\ref{fig:spectrum}).
\begin{figure}[ht]
\centering\figfile{spectrum}
\caption{Spectrum of $(\hat H-E_0,\hat{\PPv})$ in the quantum KG theory. The
vacuum sits at the origin and the one-particle states $\ket{\pp}$ on the mass
shell $(E-E_0)^2-\PPv^2=m^2$. The two-particle states $\ket{\pp_1\pp_2}$ fill
the region $(E-E_0)^2-\PPv^2\ge4m^2$, the three-particle states the region
$(E-E_0)^2-\PPv^2\ge9m^2$, and so on.}
\label{fig:spectrum}
\end{figure}
Correspondingly, the spectrum of the operator
$\hat M^2=(\hat H-E_0)^2-\hat{\PPv}^2$ has two discrete components, $0$ and
$m^2$, and continuous part from $4m^2$ to infinity.

\subsection*{Klein-Gordon propagator}

By their definition, the creation and annihilation operators
$\hat A^\dagger_\pp,\hat A_\pp$ are the Fourier components of the local field
operators $\hat\varphi(\xx)$ and $\hat\pi(\xx)$,
\begin{align*}
\hat\varphi(\xx)&=\int d\mu(\pp)\,\bigl[\hat A_\pp+\hat A^\dagger_{-\pp}\bigr]\,e^{i\pp\cdot\xx}\,,\\
\hat\pi(\xx)&=\int d\mu(\pp)\,i\omega_\pp\,\bigl[-\hat A_\pp+\hat A^\dagger_{-\pp}\bigr]\,e^{i\pp\cdot\xx}\,,
\end{align*}
where I have introduced special notation for the Lorentz-invariant measure on
the mass shell,
\[
d\mu(\pp)=\frac1{2\omega_\pp}\,\frac{d^3\pp}{(2\pi)^3}\,.
\]

\begin{exercise}
Check that the measure $d\mu(\pp)$ is Lorentz-invariant. An instructive way to
do that is to check the identity
\[
d\mu(\pp)=\int_{p^0>0}\frac{dp^0}{2\pi}\;2\pi\,\delta(p^\mu p_\mu-m^2)\,\frac{d^3\pp}{(2\pi)^3}\,,
\]
i.e.\ we have
\[
\int d\mu(\pp)\,F(\pp,\omega_\pp)=\int_{p^0>0}\frac{d^4p}{(2\pi)^4}\,2\pi\,\delta(p^\mu p_\mu-m^2)\,F(\pp,p^0)\,.
\]
\end{exercise}

The operators $\hat\varphi(\xx)$ and $\hat\pi(\xx)$ are the
``Schroedinger'', or the ``equal time'' field operators. As defined, they obey
the canonical commutation relations
\[
\bigl[\hat\pi(\xx),\hat\varphi(\xx')\bigr]=-i\,\delta^{(3)}(\xx-\xx')\,.
\]
As is common in quantum mechanics, it is often useful to trade the time
evolution of the states for the time evolution of the operators. This is known
as the Heisenberg picture. In our case, introduce the Heisenberg field
operator
\[
\hat\varphi(x)=\hat\varphi(\xx,t)=e^{i\hat Ht}\,\hat\varphi(\xx)\,e^{-i\hat Ht}\,,\qquad x=(\xx,t)\,.
\]
In terms of the creation and annihilation operators it is
\[
\hat\varphi(x)=\int d\mu(\pp)\,\bigl[\hat A_\pp\,f^*_\pp(x)+\hat A^\dagger_\pp\,f_\pp(x)\bigr]\,,
\]
where as before $f_\pp(x)=e^{i\omega_\pp t-i\pp\cdot\xx}$. Then
\[
\hat\varphi(\xx)=\hat\varphi(\xx,t=0)\,,\qquad \hat\pi(\xx)=\partial_t\hat\varphi(\xx,t)\big|_{t=0}\,.
\]

\subsection*{Commutators}

It is interesting and instructive to see how the canonical commutators of the
``Schroedinger'' operators extend to the Heisenberg field operators
$\hat\varphi(x)$. For this calculation, it is useful to write
\[
\hat\varphi(x)=\hat\varphi_-(x)+\hat\varphi_+(x)\,,
\]
where
\[
\hat\varphi_-(x)=\int d\mu(\pp)\,\hat A_\pp\,f^*_\pp(x)\,,\qquad
\hat\varphi_+(x)=\int d\mu(\pp)\,\hat A^\dagger_\pp\,f_\pp(x)\,.
\]
The operator $\hat\varphi_+(x)$ contains only the creation operators
$\hat A^\dagger_\pp$; it can be interpreted as an operator creating particle
at the space-time point $x$; likewise, $\hat\varphi_-(x)$ absorbs the particle
at the space-time point $x$.\footnote{A subtlety in this interpretation will
be discussed later in Lecture~6.} Since $\hat\varphi_-(x)$ obviously commutes
with $\hat\varphi_-(x')$, and the same is true for the $\hat\varphi_+$
component, we have
\[
\bigl[\hat\varphi(x),\hat\varphi(x')\bigr]=D_-(x-x')-D_+(x-x')\,,
\]
where
\begin{align*}
D_-(x-x')&=\bigl[\hat\varphi_-(x),\hat\varphi_+(x')\bigr]\,,\\
D_+(x-x')&=\bigl[\hat\varphi_-(x'),\hat\varphi_+(x)\bigr]\equiv D_-(x'-x)\,.
\end{align*}
Explicitly
\begin{align*}
D_-(\xx,t)&=\int d\mu(\pp)\,e^{-i\omega_\pp t+i\pp\cdot\xx}\,,\\
D_+(\xx,t)&=\int d\mu(\pp)\,e^{i\omega_\pp t-i\pp\cdot\xx}\,.
\end{align*}
Some properties of these integrals are worth mentioning.

i) The integrals for $D_-(x)$ and $D_+(x)$ are Lorentz-invariant, i.e.\ if two
4-vectors $x$ and $x'$ are related by a Lorentz transformation, $x'=\Lambda x$,
then $D(x)=D(x')$. Indeed, the integrals can be written in explicitly
covariant form, say
\[
D_-(x)=\int_{p^0>0}\frac{d^4p}{(2\pi)^4}\,2\pi\,\delta(p^\mu p_\mu-m^2)\,e^{-ip_\mu x^\mu}\,.
\]
Then the change of the integration variables $p=\Lambda^{-1}p'$ verifies that
$D_-(x)=D_-(\Lambda x)$. Similar representation exists for $D_+(x)$.

ii) The integrals are not absolutely convergent. It is usually convenient to
compute such integrals as limiting values of absolutely convergent integrals.
In our case it can be done by considering the integrals at complex values of
the time variable $t$. For instance, the above integral defining
$D_-(\xx,t)$ converges absolutely if $t$ has \textbf{negative imaginary part}
(remember that $\omega_\pp$ is positive), therefore this integral defines the
function $D_-(\xx,t)$ of complex variable $t$ which is analytic in the lower
half-plane of the complex $t$-plane. Then, for real $t$ we need to take the
limit $\lim_{\epsilon\to-0}D_-(\xx,t+i\epsilon)$ (this is usually written as
$D_-(\xx,t-i0)$). Likewise, the integral for $D_+$ defines the function
$D_+(\xx,t)$ of complex $t$, analytic in the upper half-plane, and for real
$t$ we take $D_+(\xx,t+i0)$.

Let us consider the case of real $t$ such that the 4-vector $x^\mu=(t,\xx)$ is
space-like,
\[
x^\mu x_\mu=t^2-\xx^2<0
\]
(in this consideration we assume that $\xx$ is real). In the complex $t$-plane
we are looking at the segment
\[
-|\xx|<\Re t<|\xx|
\]
of the real axis. It is easy to show that at this segment we have
\[
D_+(\xx,t)=D_-(\xx,t)\qquad\text{for}\quad -|\xx|<t<|\xx|\,.
\]
It suffices to verify this identity at $t=0$, since one can always transform
to the frame where $t=0$, as long as $x^\mu$ is space-like. For $t=0$ we have
\[
D_-(\xx,0)=\lim_{\epsilon\to+0}\int e^{-\omega_\pp\epsilon+i\pp\cdot\xx}\,d\mu(\pp)
\qquad\text{and}\qquad
D_+(\xx,0)=\lim_{\epsilon\to+0}\int e^{-\omega_\pp\epsilon-i\pp\cdot\xx}\,d\mu(\pp)\,,
\]
and the identity is established by the change of the integration variables
$\pp\to-\pp$. Note that this argument also shows that $D_\pm(\xx,t)$ return
real values when $t$ takes values within the above segment of the real axis.

It follows that the function $D_-(\xx,t)$, which was originally defined (by
our integral) in the lower half-plane of $t$, in fact can be analytically
continued, through the above segment of the real axis, to the upper
half-plane, and there it coincides with $D_+(\xx,t)$ (defined, again, through
its integral representation written above); see figure~\ref{fig:tplane}.
\begin{figure}[ht]
\centering\figfile{tplane}
\caption{The complex $t$-plane at fixed $\xx$. The function $D(\xx,t)$ is
analytic away from the two cuts $t\le-|\xx|$ and $t\ge|\xx|$ along the real
axis. Its values in the upper half-plane are given by the integral for
$D_+(\xx,t)$, and in the lower half-plane by the integral for $D_-(\xx,t)$.
The two continue into each other through the segment $-|\xx|<t<|\xx|$.}
\label{fig:tplane}
\end{figure}

In other words, in fact there is a single function of complex $t$ (I will
denote it as $D(\xx,t)$), analytic in both the upper and lower half-planes,
including the segment $-|\xx|<t<|\xx|$ of the real axis, and $D_+(\xx,t)$
($D_-(\xx,t)$) is given by its values in the upper (lower) half-plane. In
particular, for real $t$
\[
D_+(\xx,t)=D(\xx,t+i0)\,,\qquad D_-(\xx,t)=D(\xx,t-i0)\,.
\]
As we have seen, these two limits coincide when $-|\xx|<t<|\xx|$, but outside
this segment they don't. That means the function $D(\xx,t)$ has two branch
cuts in the complex $t$-plane, from $-\infty$ to $-|\xx|$, and from $|\xx|$ to
$+\infty$.

It is not difficult to identify the geometric meaning of the branching points
$t=\pm|\xx|$. They represent points on the future and past components of the
light cone associated with the space-time point $0$, that is the geometric
place of all points which can be connected to $0$ by light rays
(figure~\ref{fig:lightcone}). At fixed $\xx$ the real $t$ line intersects the
light cone at two points $\pm|\xx|$. When $t$ is regarded as the complex
variable, the intersection points become exactly the branching points in the
complex $t$ plane. At real $t$, when crossing the light cone one needs to
choose the branch of the analytic function (similar to choosing a branch of
$\sqrt z$ at negative $z$). The functions $D_+$ and $D_-$ are the two branches
corresponding to two ways of going around the branching points in the complex
$t$-plane.
\begin{figure}[ht]
\centering\figfile{lightcone}
\caption{The light cone of the origin. At fixed $\xx$ the line of real $t$
(thick) meets the future and past sheets of the cone at $t=|\xx|$ and
$t=-|\xx|$. Between these points the separation from the origin is
space-like.}
\label{fig:lightcone}
\end{figure}

Let us come back to the commutator
$[\hat\varphi(x),\hat\varphi(0)]=D_-(x)-D_+(x)$. Here I have set $x'=0$; in
general case $x$ has to be replaced by the difference $x-x'$. We see that the
commutator is equal to the discontinuity
\[
\bigl[\hat\varphi(\xx,t),\hat\varphi(\mathbf0,0)\bigr]=D(\xx,t-i0)-D(\xx,t+i0)\,.
\]
When the separation $x=(t,\xx)$ is space-like the discontinuity vanishes, and
we conclude that
\[
\bigl[\hat\varphi(x),\hat\varphi(x')\bigr]=0\qquad\text{for all space-like separations}\quad x-x'\,.
\]
This is important conclusion. Recall that in quantum mechanics commutator of
two operators allows one to determine how, in a given state, the measurement
of one observable can affect results of the measurements of another
observable. Our calculation shows that the measurement at some space-time
point $x$ can not affect measurements at the point $x'$ as long as the
separation $x-x'$ is space-like. It expresses the \textbf{causality} of the
quantum field theory. This is general requirement of quantum field theory
called the \textbf{local commutativity}. It states that for any local fields
$\mathcal O_1(x)$, $\mathcal O_2(x')$ we must have
\[
\bigl[\hat{\mathcal O}_1(x),\hat{\mathcal O}_2(x')\bigr]=0\qquad\text{for all space-like separations}\quad x-x'\,.
\]
The discontinuities across the branch cuts $|t|>|\xx|$ do not vanish, so when
the separation is time-like the commutator takes non-zero values. We see that
the commutator has its support inside the light cone.

Consider now the expectation value $\bra0\hat\varphi(x)\,\hat\varphi(x')\ket0$
of the product of two field operators. In the KG theory this expectation
value can be expressed through the same analytic function $D(\xx,t)$. Indeed,
since $\hat\varphi_-(x)\ket0=0$ and $\bra0\hat\varphi_+(x)=0$, we have
\[
\bra0\hat\varphi(x)\,\hat\varphi(x')\ket0=\bra0\hat\varphi_-(x)\,\hat\varphi_+(x')\ket0
=D_-(x-x')=D(\xx-\xx',t-t'-i0)\,.
\]
If we interchange the positions of the operators, by similar calculation
\[
\bra0\hat\varphi(x')\,\hat\varphi(x)\ket0=D_+(x-x')=D(\xx-\xx',t-t'+i0)\,.
\]
Thus, these expectation values are described by the function $D(\xx,t)$ taken
at the upper or lower edges of the light-cone branch cuts
(figure~\ref{fig:tcontours}).
\begin{figure}[ht]
\centering\figfile{tcontours}
\caption{Contours in the complex $t$-plane. Along $C_+$, just above both cuts,
$D$ gives $\bra0\hat\varphi(\mathbf0,0)\,\hat\varphi(\xx,t)\ket0$; along
$C_-$, just below both cuts, it gives
$\bra0\hat\varphi(\xx,t)\,\hat\varphi(\mathbf0,0)\ket0$. The Feynman contour
$C_F$ runs above the left cut and below the right cut.}
\label{fig:tcontours}
\end{figure}

\subsection*{Feynman Propagator}

Important role in our future studies will belong to the \textbf{time ordered}
expectation values
\[
D_F(x-x')=D_F(\xx-\xx',t-t')=\bra0\Tprod\bigl(\hat\varphi(\xx,t)\,\hat\varphi(\xx',t')\bigr)\ket0\,,
\]
where the symbol $\Tprod$ signifies the time ordered product of the operators,
\[
\Tprod\bigl(\hat\varphi(\xx,t)\,\hat\varphi(\xx',t')\bigr)=
\begin{cases}
\hat\varphi(\xx,t)\,\hat\varphi(\xx',t') & \text{for}\quad t>t'\\
\hat\varphi(\xx',t')\,\hat\varphi(\xx,t) & \text{for}\quad t<t'\,.
\end{cases}
\]
This expectation value is also called the \textbf{Feynman propagator}. One can
visualize it in figure~\ref{fig:tcontours} by first considering complex values
of $t$ along the contour $C_F$, and then taking the limit when this contour is
brought to the real axis; $D_F$ corresponds to the values of $D$ along $C_F$.

The Feynman propagator can be written as
\[
D_F(\xx,t)=
\begin{cases}
D_-(\xx,t)=\int d\mu(\pp)\,e^{-i\omega_\pp t+i\pp\cdot\xx} & \text{for}\quad t>0\\[2pt]
D_+(\xx,t)=\int d\mu(\pp)\,e^{i\omega_\pp t-i\pp\cdot\xx} & \text{for}\quad t<0\,,
\end{cases}
\]
but the most convenient representation is given by the 4-dimensional integral
\[
D_F(\xx,t)=\int\frac{d^3\pp}{(2\pi)^3}\int_{\tilde C_F}\frac{d\omega}{2\pi}\,
\frac{i\,e^{-i\omega t+i\pp\cdot\xx}}{\omega^2-\pp^2-m^2}\,.
\]
The integrand has two poles in $\omega$, at $\omega=\pm\omega_\pp$, and the
integration contour $\tilde C_F$ for $\omega$ goes around these poles as shown
in figure~\ref{fig:omegaF}.
\begin{figure}[ht]
\centering\figfile{omegaF}
\caption{The contour $\tilde C_F$ in the complex $\omega$-plane passes below
the pole at $\omega=-\omega_\pp$ and above the pole at $\omega=\omega_\pp$.}
\label{fig:omegaF}
\end{figure}

If $t>0$ the integrand decays exponentially into the lower half-plane of the
variable $\omega$, and one can perform the integration over this variable by
closing the contour around $\omega=\omega_\pp$. The integral reduces to the
residue at this pole, which yields $D_-(\xx,t)$. Similarly, when $t<0$ the
integration contour can be closed around the pole $\omega=-\omega_\pp$, and the
residue calculation yields $D_+(\xx,t)$.

This integration prescription is usually expressed as
\[
D_F(x)=\int\frac{d\omega}{2\pi}\,\frac{d^3\pp}{(2\pi)^3}\,\frac{i\,e^{-i\omega t+i\pp\cdot\xx}}{\omega^2-\pp^2-m^2+i0}
=\int\frac{d^4p}{(2\pi)^4}\,\frac{i\,e^{-ip_\mu x^\mu}}{p_\mu p^\mu-m^2+i0}\,.
\]
This form of the Feynman propagator is explicitly covariant. It also shows that
$D_F(x)$ is the Green's function of the KG equation,
\[
\bigl(\partial_\mu\partial^\mu+m^2\bigr)\,D_F(x)=-i\,\delta^{(4)}(x)\,.
\]
Note that the functions $D_\pm(x)$ satisfy homogeneous KG equation
\[
\bigl(\partial_\mu\partial^\mu+m^2\bigr)\,D_\pm(x)=0\,,
\]
since in the KG theory the Heisenberg operator $\hat\varphi(x)$ itself
satisfies homogeneous KG equation.

\begin{exercise}
The Feynman propagator is defined as
\[
D_F(\xx,t)=\theta(t)\,\bra0\hat\varphi(\xx,t)\,\hat\varphi(\mathbf0,0)\ket0+\theta(-t)\,\bra0\hat\varphi(\mathbf0,0)\,\hat\varphi(\xx,t)\ket0\,.
\]
Applying the differential operator $\partial_\mu\partial^\mu+m^2$ and using the
equal-time commutation relations
\[
\bigl[\hat\pi(\xx,t),\hat\varphi(\xx',t)\bigr]=-i\,\delta^{(3)}(\xx-\xx')\,,
\]
confirm that
\[
\bigl(\partial_\mu\partial^\mu+m^2\bigr)\,D_F(x)=-i\,\delta^{(4)}(x)\,.
\]
\end{exercise}

Finally, let us consider the ``euclidean'', or ``imaginary time'' correlation
function. It is defined in terms of the function $D(\xx,t)$ taken at pure
imaginary values of the time variable. Namely, let
\[
t=-i\,x_4\,,
\]
with real $x_4$. Then
\[
D_E(\xx,x_4)=D(\xx,-ix_4)\,,
\]
i.e.\ $D_E$ is the function $D$ with $t$ taken along the contour $C_E$ in the
complex $t$-plane (figure~\ref{fig:tE}).
\begin{figure}[ht]
\centering\figfile{tE}
\caption{The Euclidean contour $C_E$ runs along the imaginary $t$ axis,
between the two light-cone cuts.}
\label{fig:tE}
\end{figure}
Note that this contour can be obtained from $C_F$ by $90^\circ$ rotation.

As is well known, by taking pure imaginary values of $t$ one converts the
Minkowski space-time into 4-dimensional Euclidean space, since
\[
-d\tau^2\to ds^2=dx_4^2+d\xx^2\,,
\]
and the Lorentz symmetry becomes the orthogonal group $O(4)$ of rotations of
the Euclidean 4-vectors
\[
x_E=(\xx,x_4)\,.
\]

The function $D_E(x_E)$ is called the Euclidean correlation function of the KG
theory. We will discuss its significance later. For now, let us observe that
since $D_E$ is related to $D_F$ simply by rotating the time contour, it
satisfies the Euclidean-space version of the inhomogeneous KG equation
\[
\Bigl(m^2-\frac{\partial^2}{\partial x_4^2}-\nabla^2\Bigr)\,D_E(\xx,x_4)=\delta(x_4)\,\delta^{(3)}(\xx)\,,
\]
or in symmetric form
\[
\bigl(m^2-\Delta_E\bigr)\,D_E(x_E)=\delta^{(4)}(x_E)\,,
\]
where $\Delta_E$ stands for the 4-dimensional Laplacian. This equation is
easily solved by Fourier transformation
\[
D_E(x_E)=\int\frac{d^4p_E}{(2\pi)^4}\,\frac{e^{ip_Ex_E}}{p_E^2+m^2}\,,
\]
where the integration is over the 4-dimensional Euclidean momentum space,
$p_E=(\pp,p_4)$, and $p_E^2=\pp^2+p_4^2$. Note that now the integrand does not
have singularities in the integration domain. Also, unlike the Minkowski space
version of this equation, the solution is unique if we demand that $D_E(x_E)$
decays at the Euclidean infinity $x_E\to\infty$. By introducing auxiliary
integration, one can write this function as
\begin{equation}\label{eq:Schwinger}
D_E(x_E)=\int_0^\infty d\tau\int\frac{d^4p_E}{(2\pi)^4}\,e^{-\tau(p_E^2+m^2)}\,e^{ip_Ex_E}\,.
\end{equation}
This form is known as the Schwinger's proper-time representation.

Clearly, the above Euclidean momentum space integral for $D_E(x_E)$ is related
to the covariant 4-momentum integral for the Feynman propagator. The
correlation function $D_E(\xx,x_4)$ is the analytic continuation of
$D_F(\xx,t)$ obtained by rotating the time contour $C_F$ to $C_E$.
Correspondingly, the integration contour $\tilde C_F$ in the complex
$\omega$-plane is rotated by $90^\circ$ to the contour $\tilde C_E$
(figure~\ref{fig:omegaE}).
\begin{figure}[ht]
\centering\figfile{omegaE}
\caption{Wick rotation in the $\omega$-plane. The contour $\tilde C_F$ can be
turned counterclockwise by $90^\circ$ into $\tilde C_E$, along the imaginary
axis, without crossing the poles at $\omega=\pm\omega_\pp$.}
\label{fig:omegaE}
\end{figure}

This example suggests that some analytic characteristics of the KG theory,
and perhaps more general QFT, may take simpler and more symmetric form when
continued to pure imaginary time. This idea can be refined with the help of
the path integral formulation of quantum mechanics, which we consider next.

\begin{hwproblem}{7}\label{hw:7}
\begin{hwparts}
\item Show that the Euclidean Klein-Gordon propagator
\[
D_E(x_E):=\int\frac{d^4p}{(2\pi)^4}\,\frac{e^{ipx_E}}{p^2+m^2}
\]
behaves as
\[
D_E(x_E)\to\Const\,\frac{m^{1/2}}{R^{3/2}}\,e^{-mR}\qquad\text{for}\quad mR\to\infty\,,
\]
where $R$ stands for $|x_E|$. Evaluate the $\Const$.
\item Using the Schwinger's representation, find $D_E$ explicitly in terms of
the Bessel's functions. (Consult e.g.\ A.~Erd\'elyi, H.~Bateman, ``Higher
Transcendental Functions'', Volume~II,
\url{https://authors.library.caltech.edu/43491/}).
\item Analyze the singularity at $x_E^2=0$. For $x_E=(\xx,x_4)$, consider
analytic continuation to complex $x_4=it$. Describe analytic structure of
$D(\xx,t)$ as the function of complex $t$. Find the discontinuity across the
branch cuts associated with the future and past components of the light cone
$t^2=\xx^2$.
\end{hwparts}
\end{hwproblem}
